\section*{Capítulo \ref{ch:b1:alcohol-activation} --- Álcoois: ativar o grupo OH}

\begin{solution}{exo:b1:alcohol-activation:1}
\ce{2CH3CH2OH + 2Na -> 2CH3CH2O- + 2Na+ + H2};
\ce{CH3CH2OH + NaH -> CH3CH2O- + Na+ + H2};
\ce{CH3CH2OH + OH- <=> CH3CH2O- + H2O}, $K = 10^{14.0 - 15.9} = 10^{-1.9}$:
parcial.
\end{solution}

\begin{solution}{exo:b1:alcohol-activation:2}
1-Metoxipropano; metoxietano; etoxibenzeno (o fenóxido, base mais fraca
que um \omterm{def:b1:alcohol-activation:alkoxide}{alcóxido}, continua sendo um bom \omterm{def:b1:electronic-effects:nucleophile}{nucleófilo}).
\end{solution}

\begin{solution}{exo:b1:alcohol-activation:3}
A tosilação em piridina dá \ce{CH3CH2CH2OTs}; depois
\[ \ce{CH3CH2CH2OTs + I- -> CH3CH2CH2I + TsO-}. \] Diretamente, o iodeto teria de
expulsar \ce{OH-}: nenhuma reação. O \omterm{def:b1:alcohol-activation:sulfonate}{tosilato} tem um excelente \omterm{def:b1:electronic-effects:nucleophile}{grupo de saída},
e as condições não são nem ácidas nem fortemente básicas.
\end{solution}

\begin{solution}{exo:b1:alcohol-activation:4}
Butan-2-ol: but-1-eno, \cip{Z}- e \cip{E}-but-2-eno; majoritário:
\cip{E}-but-2-eno. 2-Metilpentan-2-ol: 2-metilpent-2-eno (majoritário,
trissubstituído) e 2-metilpent-1-eno.
\end{solution}

\begin{solution}{exo:b1:alcohol-activation:5}
\ce{(CH3)2CHCH2-O-CH2CH3}: 2-metilpropan-1-olato de sódio com bromoetano,
ou etóxido de sódio com 1-bromo-2-metilpropano. Os dois haletos são
primários, mas o 1-bromo-2-metilpropano é ramificado ao lado do carbono
que reage, o que torna a SN2 mais lenta e deixa a \omterm{def:b1:predominant-reaction:strong-acid}{base forte} eliminar
(2-metilpropeno). A primeira escolha é melhor.
\end{solution}

\begin{solution}{exo:b1:alcohol-activation:6}
O \omterm{def:b1:alcohol-activation:sulfonate}{tosilato} conserva a \omterm{def:b1:stereochemistry-in-depth:isomers}{configuração} \cip{R}; a SN2 com o etóxido a inverte:
\cip{S}-2-etoxioctano (OEt é o primeiro na ordem, como era OTs). Com ácido
sulfúrico e calor: eliminação em octenos (principalmente o
\cip{E}-oct-2-eno), via um \omterm{def:b1:electronic-effects:intermediates}{carbocátion} secundário, com perda da estereoquímica.
\end{solution}

\begin{solution}{exo:b1:alcohol-activation:7}
Protonação do OH; perda de água até o \omterm{def:b1:electronic-effects:intermediates}{carbocátion} terciário (lenta);
captura por \ce{Cl-}: 2-cloro-2-metilpropano. Rápida porque o \omterm{def:b1:electronic-effects:intermediates}{carbocátion}
terciário é estável e o cloreto, insolúvel no ácido, se separa.
\end{solution}

\begin{solution}{exo:b1:alcohol-activation:8}
\ce{CH3CH2OH + H+ -> CH3CH2OH2+}; um segundo etanol ataca o carbono por
trás enquanto a água sai (SN2): \ce{(CH3CH2)2OH+}, que perde um próton:
etoxietano. Concorrente: uma base (etanol, \ce{HSO4-}) retira um próton $\beta$
de \ce{CH3CH2OH2+} enquanto a água sai (E2): eteno.
\end{solution}

\begin{solution}{exo:b1:alcohol-activation:9}
2-Metilpropan-2-ol (terciário), dando 2-cloro-2-metilpropano, insolúvel.
\end{solution}

\begin{solution}{exo:b1:alcohol-activation:10}
A protonação e a perda de água dão um \omterm{def:b1:electronic-effects:intermediates}{carbocátion} secundário vizinho de
um carbono quaternário. Um grupo metila desse carbono migra com seu \omterm{def:b1:lewis-resonance-vsepr:lewis}{par
ligante} para o carbono catiônico: a carga positiva fica agora em um
\omterm{def:b1:electronic-effects:carbon-class}{carbono terciário}, mais estável. A perda de um próton de um carbono
vizinho dá o 2,3-dimetilbut-2-eno tetrassubstituído, o alceno mais estável.
\end{solution}

\begin{solution}{exo:b1:alcohol-activation:11}
A partir do \cip{R}-butan-2-ol: \omterm{def:b1:alcohol-activation:sulfonate}{tosilato} (\cip{R}), depois metóxido de sódio
(SN2, inversão): \cip{S}-2-metoxibutano. A partir do \cip{S}-butan-2-ol:
prepare o \omterm{def:b1:alcohol-activation:alkoxide}{alcóxido} (NaH) e adicione iodometano: a ligação C--O
estereogênica não é tocada, e de novo se obtém o \cip{S}-2-metoxibutano.
\end{solution}

\begin{solution}{exo:b1:alcohol-activation:12}
$Q = [\text{éter}][\ce{H2O}]/[\ce{EtOH}]^2$. Remover a água à medida que
ela se forma mantém $[\ce{H2O}]$ e $Q$ pequenos: $Q < K$ e a reação
continua formando éter até o álcool ser consumido.
\end{solution}

\begin{solution}{pb:b1:alcohol-activation:1}
\textbf{1.} Terc-butóxido de sódio com um haleto de metila; metóxido de
sódio com um haleto de terc-butila.
\textbf{2.} O segundo: o metóxido, uma \omterm{def:b1:predominant-reaction:strong-acid}{base forte}, provoca eliminação no
haleto terciário, \ce{(CH3)3CBr + CH3O- -> (CH3)2C=CH2 + CH3OH + Br-}.
\textbf{3.} \ce{(CH3)3CO- + CH3I -> (CH3)3COCH3 + I-}: o oxigênio do
\omterm{def:b1:alcohol-activation:alkoxide}{alcóxido} ataca o carbono do metila por trás enquanto o iodeto sai (SN2).
\textbf{4.} Com sódio (ou hidreto de sódio): \ce{2(CH3)3COH + 2Na ->
2(CH3)3CO- + 2Na+ + H2}. O hidróxido é uma \omterm{def:b1:predominant-reaction:strong-acid}{base fraca} demais: $K = 10^{14.0 -
19.2} = 10^{-5.2}$.
\textbf{5.} A água protonaria o \omterm{def:b1:alcohol-activation:alkoxide}{alcóxido} e atacaria o haleto.
\textbf{6.} Não: nenhum \omterm{def:b1:stereochemistry-in-depth:stereocentre}{centro estereogênico}.
\textbf{7.} \cip{R}-butan-2-ol + TsCl (piridina) $\to$ \omterm{def:b1:alcohol-activation:sulfonate}{tosilato} de
\cip{R}-butan-2-ila: retenção.
\textbf{8.} \cip{S}-2-metoxibutano, por inversão SN2.
\textbf{9.} \cip{S}-butan-2-ol.
\textbf{10.} \cip{R}-2-metoxibutano: a ligação C--O do carbono
estereogênico é conservada; só a ligação O--C(metila) se forma.
\textbf{11.} A E2 (o metóxido é uma \omterm{def:b1:predominant-reaction:strong-acid}{base forte} e o carbono é secundário),
que dá butenos; limitada pela base pequena, por um excelente \omterm{def:b1:electronic-effects:nucleophile}{grupo de
saída} e por uma temperatura moderada.
\textbf{12.} O álcool protonado se ioniza em um \omterm{def:b1:electronic-effects:intermediates}{carbocátion} secundário
plano, capturado pelas duas faces pelo metanol: éter racêmico,
acompanhado de butenos.
\textbf{13.} Protonação do OH, perda de água (lenta) até o \omterm{def:b1:electronic-effects:intermediates}{carbocátion}
terciário, perda de um próton $\beta$.
\textbf{14.} 2-Metilbut-2-eno (majoritário, trissubstituído) e
2-metilbut-1-eno.
\textbf{15.} Um \omterm{def:b1:electronic-effects:intermediates}{carbocátion} terciário é impedido e perde um próton muito
mais facilmente do que outra molécula de álcool consegue atacá-lo.
\textbf{16.} Para remover um produto: com $Q$ mantido pequeno a reação
continua, e o alceno não fica em contato com o ácido.
\textbf{17.} Uma primeira barreira alta (perda de água, determinante) até
o \omterm{def:b1:electronic-effects:intermediates}{carbocátion}, depois uma barreira pequena até o alceno.
\textbf{18.} Não: um \omterm{def:b1:electronic-effects:intermediates}{carbocátion} primário é instável demais; os álcoois
primários eliminam por E2 sobre o álcool protonado, a temperatura mais alta.
\textbf{19.} $10.0/74.0 = \qty{0.135}{mol}$.
\textbf{20.} $1.10 \times 0.135 \times 141.9 = \qty{21.1}{g}$.
\textbf{21.} $0.135 \times 88.0 = \qty{11.9}{g}$.
\textbf{22.} $0.75 \times 11.9 =$ \textbf{\qty{8.9}{g} de MTBE}.
\end{solution}
